% Abhakseite „Das kann ich“ – Zone nur im Gesamt (\mitzone)
%% TODO Ich-kann-Sätze: Platzhalter wie in den Aufgabendateien
\begin{abhakseite}
\ifdefined\mitzone
\abhakgruppe{Kennst du schon}
\abhak{1}{Fläche und Umfang unterscheiden; Einheiten cm, cm², m, m²}
\abhak{2}{Mit Dezimalzahlen multiplizieren und dividieren, Taschenrechner mit π, Ergebnis runden (Näherungswert mal Kommazahl)}
\abhak{3}{Formel nach einer Größe umstellen (u = π · d → d = u : π)}
\abhak{4}{Quadrieren und Wurzelziehen (4,5²; √20,25)}
\abhak{5}{Winkel messen und zeichnen, Vollwinkel}
\abhak{6}{Anteil als Bruch und Prozentsatz bilden (Mittelpunktswinkel zum Vollwinkel als Prozent)}
\abhak{7}{Fläche und Umfang unterscheiden; Einheiten cm, cm², m, m² – Fehler finden}
\abhak{8}{Fläche und Umfang unterscheiden; Einheiten cm, cm², m, m² – selbst rechnen}
\fi
\abhakgruppe{Einheit 1 · Kreisumfang}
\abhak{9}{Rand oder Fläche?}
\abhak{10}{Umfang}
\abhak{11}{Radius, Durchmesser, Mittelpunkt in einer Figur benennen und einzeichnen}
\abhak{12}{Kreis mit gegebenem Radius oder Durchmesser zeichnen}
\abhak{13}{Umfang messen und u : d bilden (Entdeckung von π)}
\abhak{14}{Umfang – Fehler finden, Begründen, Anwendung, Darstellungswechsel}
\abhakgruppe{Einheit 2 · Kreisfläche}
\abhak{15}{Fläche}
\abhak{16}{Term zu Figur}
\abhak{17}{Sachaufgabe (Abwurfring, Pizza, Deckel, Grundfläche eines Kegels oder Zylinders)}
\abhak{18}{Fläche oder Umfang: die passende Formel wählen}
\abhak{19}{Fläche – Fehler finden, Begründen, Anwendung, Darstellungswechsel}
\abhakgruppe{Einheit 3 · Kreisteile}
\abhak{20}{Kreisteile}
\abhak{21}{Sachaufgabe (Rasensprenger, Tortenstück, Sektor im Kreisdiagramm)}
\abhak{22}{Kreisteile – Fehler finden, Begründen, Anwendung, Darstellungswechsel}
\end{abhakseite}
